% --- Chapter: References and Pointer Semantics ---
\section{References and Pointer Semantics}

CRISP provides shared mutable references through a pointer-like abstraction. Unlike C pointers (raw memory addresses) or C++ references (non-nullable aliases), CRISP references use reference-counted interior mutability — safe, shared ownership of a mutable value.

\subsection{Creating and Dereferencing}

A reference is created with the \texttt{\textbackslash} operator and dereferenced with \texttt{\textasciicircum}:

\begin{tcolorbox}[colback=codebg, colframe=darkpurple, colbacktitle=darkpurple, title=\textbf{\textcolor{codegold}{Creating and Dereferencing}}, fonttitle=\bfseries, sharp corners]
\begin{lstlisting}
# Create a reference to a value
let x = 42;
let r = \x;

# Dereference to read
say ^r;       # 42

# Dereference to write (modify through reference)
^r = 100;
say x;        # 100 — the original was mutated!

# Multiple references to the same value
let a = 10;
let ref1 = \a;
let ref2 = \a;

^ref1 = 20;
say ^ref2;    # 20 — both references point to the same cell
say a;        # 20
\end{lstlisting}
\end{tcolorbox}

\subsection{The Reference Sigil (\texttt{\&})}

Variables holding references can be declared with the \texttt{\&} sigil for clarity:

\begin{tcolorbox}[colback=codebg, colframe=darkpurple, colbacktitle=darkpurple, title=\textbf{\textcolor{codegold}{Reference Sigil}}, fonttitle=\bfseries, sharp corners]
\begin{lstlisting}
# With sigil (explicit — "this is a reference")
let &counter = \0;
^counter = ^counter + 1;
say ^counter;   # 1

# Without sigil (implicit — works the same)
let counter = \0;
^counter = ^counter + 1;
say ^counter;   # 1

# The sigil is a naming convention, not a type constraint
let &r = \0;
r = \100;       # reassign to reference a different value
say ^r;         # 100
\end{lstlisting}
\end{tcolorbox}

\subsection{References to Collections}

References to arrays and hashes enable shared mutable access to data structures without copying:

\begin{tcolorbox}[colback=codebg, colframe=darkpurple, colbacktitle=darkpurple, title=\textbf{\textcolor{codegold}{References to Collections}}, fonttitle=\bfseries, sharp corners]
\begin{lstlisting}
# Reference to an array
let @data = [1, 2, 3];
let &r = \@data;

push(^r, 4);
push(^r, 5);

say @data;      # [1, 2, 3, 4, 5]
say ^r;         # [1, 2, 3, 4, 5]

# Reference to a hash
let %config = { debug => false };
let &cfgref = \%config;

# Modify through reference
let h = ^cfgref;
h["debug"] = true;

say %config["debug"];  # true
\end{lstlisting}
\end{tcolorbox}

\subsection{References as Function Arguments}

Passing a reference to a function enables the function to mutate the caller's data — a form of pass-by-reference:

\begin{tcolorbox}[colback=codebg, colframe=darkpurple, colbacktitle=darkpurple, title=\textbf{\textcolor{codegold}{References as Arguments}}, fonttitle=\bfseries, sharp corners]
\begin{lstlisting}
# Function that mutates through a reference
fn increment(&ref) {
    ^ref = ^ref + 1;
}

let x = 0;
increment(\x);
increment(\x);
increment(\x);
say x;  # 3

# Swap two values through references
fn swap(&a, &b) {
    let temp = ^a;
    ^a = ^b;
    ^b = temp;
}

let first = "hello";
let second = "world";
swap(\first, \second);
say first;   # world
say second;  # hello
\end{lstlisting}
\end{tcolorbox}

\subsection{Reference Semantics: Shallow vs. Deep}

A reference points to the \textit{same} value, not a copy. This means changes through any reference are visible through all others:

\begin{tcolorbox}[colback=codebg, colframe=darkpurple, colbacktitle=darkpurple, title=\textbf{\textcolor{codegold}{Reference Semantics}}, fonttitle=\bfseries, sharp corners]
\begin{lstlisting}
# No reference — independent copies
let a = [1, 2, 3];
let b = a;          # shallow copy of the array handle
push(b, 4);
say a;              # [1, 2, 3] — unaffected
say b;              # [1, 2, 3, 4]

# With reference — shared ownership
let c = [1, 2, 3];
let &rc = \c;
let &rc2 = \c;      # second reference to same array
push(^rc, 4);
say ^rc2;           # [1, 2, 3, 4] — visible through rc2
say c;              # [1, 2, 3, 4] — original is mutated
\end{lstlisting}
\end{tcolorbox}

\subsection{Chaining References}

References can be nested — a reference to a reference:

\begin{tcolorbox}[colback=codebg, colframe=darkpurple, colbacktitle=darkpurple, title=\textbf{\textcolor{codegold}{Chained References}}, fonttitle=\bfseries, sharp corners]
\begin{lstlisting}
let x = 42;
let r1 = \x;      # r1 → x
let r2 = \r1;     # r2 → r1 → x

# Double dereference
say ^^r2;         # 42

# Modify through chain
^^r2 = 99;
say x;            # 99
\end{lstlisting}
\end{tcolorbox}

\subsection{Use Cases}

\begin{table}[H]
\centering
\rowcolors{2}{softviolet!30}{white}
\caption{\textcolor{darkpurple}{Reference Use Cases}}
\vspace{0.3em}
\begin{tabular}{@{}llp{7cm}@{}}
\toprule
\rowcolor{softvioletdark}
\textbf{\textcolor{white}{Use Case}} & \textbf{\textcolor{white}{Example}} & \textbf{\textcolor{white}{Why References}} \\
\midrule
Out-parameters & \texttt{fn read\_config(\&out)} & Function mutates caller's variable \\
Shared state & \texttt{let \&state = \textbackslash @log;} & Multiple closures access same data \\
Avoid copies & \texttt{let \&big = \textbackslash @large\_array;} & Pass without cloning large structures \\
Event handlers & \texttt{counter.increment(\&state)} & Multiple handlers mutate shared state \\
Circular references & \texttt{let \%node = \{ next => \textbackslash other \}} & Build linked structures \\
\bottomrule
\end{tabular}
\end{table}

\subsection{Internal Mechanism}

CRISP references are implemented as \texttt{Value::Ref(Rc<RefCell<Value>>)}:

\begin{table}[H]
\centering
\rowcolors{2}{softviolet!30}{white}
\caption{\textcolor{darkpurple}{Reference Internals}}
\vspace{0.3em}
\begin{tabular}{@{}llp{7cm}@{}}
\toprule
\rowcolor{softvioletdark}
\textbf{\textcolor{white}{Component}} & \textbf{\textcolor{white}{Rust Type}} & \textbf{\textcolor{white}{Role}} \\
\midrule
\texttt{Rc} & \texttt{std::rc::Rc} & Reference counting — tracks how many references exist \\
\texttt{RefCell} & \texttt{std::cell::RefCell} & Interior mutability — allows mutation through shared reference \\
\texttt{Value} & \texttt{Value} & The pointed-to value (can be any CRISP type) \\
\bottomrule
\end{tabular}
\end{table}

\begin{tcolorbox}[colback=lightlavender, colframe=softvioletdark, title=\textbf{\textcolor{darkpurple}{Design Note: Why Rc<RefCell<Value>>?}}, fonttitle=\bfseries, sharp corners]
The combination of \texttt{Rc} and \texttt{RefCell} provides two guarantees:

\begin{itemize}
    \item \textbf{Shared ownership (\texttt{Rc}):} Multiple references can coexist. The value is deallocated only when the last reference is dropped. No dangling pointers.
    \item \textbf{Interior mutability (\texttt{RefCell}):} Mutation through a shared reference is safe because \texttt{RefCell} enforces Rust's borrow rules at runtime (not compile time). Attempting to mutate while another reference is borrowed causes a runtime panic rather than undefined behavior.
\end{itemize}

This is the same pattern used by Rust's own \texttt{Rc<RefCell<T>>} for shared mutable state in single-threaded contexts. CRISP is currently single-threaded, so \texttt{Rc} (not \texttt{Arc}) is sufficient.
\end{tcolorbox}

\subsection{Limitations and Caveats}

\begin{tcolorbox}[colback=softviolet!20, colframe=warningred, title=\textbf{\textcolor{warningred}{Caveats}}, fonttitle=\bfseries, sharp corners]
\begin{itemize}
    \item \textbf{No automatic dereferencing:} Unlike C++ references or Rust's \texttt{Deref}, CRISP requires explicit \texttt{\textasciicircum} to dereference. Writing \texttt{r = 5} when \texttt{r} is a reference reassigns the reference, not the pointed-to value.
    \item \textbf{Single-threaded only:} \texttt{Rc} is not \texttt{Send} or \texttt{Sync}. CRISP does not currently support multi-threading, so this is not a limitation in practice. Future multi-threaded versions would use \texttt{Arc}.
    \item \textbf{No weak references:} All references are strong (\texttt{Rc}, not \texttt{Weak}). This means circular references will leak memory — they prevent deallocation.
    \item \textbf{Runtime borrow checking:} \texttt{RefCell} panics at runtime if you violate borrow rules (e.g., mutating while iterating). This is safer than C's undefined behavior but less safe than Rust's compile-time borrow checking.
\end{itemize}
\end{tcolorbox}

\subsection{Comparison with Other Languages}

\begin{table}[H]
\centering
\rowcolors{2}{softviolet!30}{white}
\caption{\textcolor{darkpurple}{Reference/Pointer Models Across Languages}}
\vspace{0.3em}
\begin{tabular}{@{}llllp{5cm}@{}}
\toprule
\rowcolor{softvioletdark}
\textbf{\textcolor{white}{Language}} & \textbf{\textcolor{white}{Create}} & \textbf{\textcolor{white}{Access}} & \textbf{\textcolor{white}{Safety}} & \textbf{\textcolor{white}{Notes}} \\
\midrule
CRISP & \texttt{\textbackslash x} & \texttt{\textasciicircum r} & Runtime (RefCell) & Shared ownership, interior mutability \\
C & \texttt{\&x} & \texttt{*p} & None & Raw pointer, UB on misuse \\
C++ & \texttt{\&x} & Implicit & None & Non-nullable alias \\
Rust & \texttt{\&x} / \texttt{\&mut x} & Implicit & Compile-time & Borrow checker \\
Perl & \texttt{\textbackslash \$x} & \texttt{\$\$r} / \texttt{\$r->} & Runtime (refcount) & Garbage collected \\
Python & Implicit & Implicit & Runtime (GC) & Everything is a reference \\
\bottomrule
\end{tabular}
\end{table}

CRISP's model sits between Perl (explicit \textbackslash/\$ deref, reference-counted) and Rust (safe interior mutability pattern), with the explicitness of C pointers but the safety of runtime borrow checking.
